\scanpage{135}
% archive foliation: 134
sur $V$ sont en correspondance biunivoque avec l'ens $\Sigma\setminus\Sigma_{\mathbb R}$ des points non réels de $\Sigma$.

Notons que l'ens $\Sigma\setminus\Sigma_{\mathbb R}$ a deux composantes connexes \inserted{$\Sigma_\omega^*$}, qui correspondent (par Stokes) aux deux orientations de $\Sigma_{\mathbb R}$, ou encore aux deux orientations de $V$.

La composante connexe d'un $z\in\Sigma\setminus\Sigma_{\mathbb R}$ correspond à l'orientation de $V$ définie de façon évidente par la structure complexe envisagée sur $V$ (définie par $z$).
\margin{C'est faux, cf calcul plus bas}

Revenant à la description transcendante des courbes elliptiques, on voit qu'\inserted{la donnée d'}une telle courbe $E$ revient à la donnée

a) D'un $\mathbb Z$-module libre de rang $2$ $H$

b) D'une \deleted{droite} orientation $\omega$ de $V=\unsure{\Pi|H}\otimes_{\mathbb Z}\mathbb R$, ou ce qui revient au $=$, d'une des deux bases de $\bigwedge_{\mathbb Z}^{2}H$ (i.e. d'une restriction du groupe structural $\mathrm{Gl}(2,\mathbb Z)$ de $H$ à $\mathrm{Sl}(2,\mathbb Z)$), \inserted{\deleted{soit} notée aussi $\omega$}
\margin{On a $H=\pi_1(E,0_E)$\\
$V=H\otimes_{\mathbb Z}\mathbb R\simeq t_E$}

c) D'un point $z\in\Sigma_\omega(V)$,

Une \uline{rigidification de Torelli} de la courbe elliptique $E$ est la donnée d'un isomorphisme
